\section{Detailed Gaussian calculations for forward diffusion}
This section proves standard identities used in the diffusion explanation. It supplies no new research claim. Coordinates are Euclidean latent coordinates; all covariance matrices below are scalar multiples of the identity.

\begin{proposition}[Forward marginal and conditional posterior]
Let $0<\beta_t<1$, $\alpha_t=1-\beta_t$, $\bar\alpha_t=\prod_{j=1}^t\alpha_j$ and
\begin{equation}
 X_t=\sqrt{\alpha_t}X_{t-1}+\sqrt{\beta_t}Z_t,\qquad Z_t\overset{\rm iid}{\sim}\mathcal N(0,I).
\end{equation}
Conditionally on $X_0=x_0$, the marginal law is
\begin{equation}
 X_t\mid x_0\sim\mathcal N(\sqrt{\bar\alpha_t}x_0,(1-\bar\alpha_t)I).
\end{equation}
For $t>1$, the conditional law of $X_{t-1}$ given $(x_t,x_0)$ has covariance $\widetilde\beta_tI$ and mean $\widetilde\mu_t$, where
\begin{align}
 \widetilde\beta_t&=\frac{\beta_t(1-\bar\alpha_{t-1})}{1-\bar\alpha_t},\\
 \widetilde\mu_t&=
 \frac{\sqrt{\bar\alpha_{t-1}}\beta_t}{1-\bar\alpha_t}x_0+
 \frac{\sqrt{\alpha_t}(1-\bar\alpha_{t-1})}{1-\bar\alpha_t}x_t.
\end{align}
\end{proposition}
\begin{proof}
The first-step conditional mean is $\sqrt{\alpha_1}x_0$ and its conditional covariance is $\beta_1I$. Suppose the marginal formula holds at $t-1$. Independence of $Z_t$ from the past gives
\begin{align}
 \E[X_t\mid x_0]
 &=\sqrt{\alpha_t}\sqrt{\bar\alpha_{t-1}}x_0
 =\sqrt{\bar\alpha_t}x_0,\\
 \operatorname{Cov}(X_t\mid x_0)
 &=\alpha_t(1-\bar\alpha_{t-1})I+\beta_tI\\
 &=[\alpha_t+\beta_t-\alpha_t\bar\alpha_{t-1}]I
 =(1-\bar\alpha_t)I.
\end{align}
A linear combination of independent Gaussians is Gaussian, so induction proves the marginal law.

For the posterior, Bayes' rule gives, up to a normalizing constant independent of $u=X_{t-1}$,
\begin{align}
 \log q(u\mid x_t,x_0)
 &= -\frac{\|x_t-\sqrt{\alpha_t}u\|^2}{2\beta_t}
    -\frac{\|u-\sqrt{\bar\alpha_{t-1}}x_0\|^2}{2(1-\bar\alpha_{t-1})}
    +\mathrm{const}\\
 &= -\tfrac12\left[a\,u^\top u-2u^\top b\right]+\mathrm{const},
\end{align}
where
\begin{align}
 a&=\frac{\alpha_t}{\beta_t}+\frac{1}{1-\bar\alpha_{t-1}}
 =\frac{\alpha_t(1-\bar\alpha_{t-1})+\beta_t}{\beta_t(1-\bar\alpha_{t-1})}\\
 &=\frac{1-\bar\alpha_t}{\beta_t(1-\bar\alpha_{t-1})},\\
 b&=\frac{\sqrt{\alpha_t}}{\beta_t}x_t+
      \frac{\sqrt{\bar\alpha_{t-1}}}{1-\bar\alpha_{t-1}}x_0.
\end{align}
Since $a>0$, completing the square yields
\begin{equation}
 a\,u^\top u-2u^\top b
 =a\|u-a^{-1}b\|^2-a^{-1}\|b\|^2.
\end{equation}
Thus the covariance is $a^{-1}I=\widetilde\beta_tI$, and the mean is
\begin{align}
 a^{-1}b
 &=\frac{\beta_t(1-\bar\alpha_{t-1})}{1-\bar\alpha_t}
 \left[\frac{\sqrt{\alpha_t}}{\beta_t}x_t+
       \frac{\sqrt{\bar\alpha_{t-1}}}{1-\bar\alpha_{t-1}}x_0\right]\\
 &=\frac{\sqrt{\alpha_t}(1-\bar\alpha_{t-1})}{1-\bar\alpha_t}x_t+
   \frac{\sqrt{\bar\alpha_{t-1}}\beta_t}{1-\bar\alpha_t}x_0.
\end{align}
This is the stated posterior mean. At $t=1$, conditioning on $X_0=x_0$ fixes the previous state exactly; the nondegenerate Gaussian posterior derivation is replaced by a point mass.
\end{proof}
